\section{Transiciones y señal de programa final}

Las transiciones controlan cómo se realiza el cambio de una fuente o composite a otro
en el flujo de programa. Voctomix 2.0 ofrece dos modos de transición que el realizador
activa desde la barra de herramientas de la interfaz gráfica.

\subsection{Cut (corte instantáneo)}

El botón \textbf{Cut} realiza una transición instantánea: en el siguiente fotograma
tras la pulsación, la fuente de programa pasa a ser la fuente de preview. No existe
ningún efecto intermedio. Es la transición estándar en producción broadcast: el
``corte'' es la forma más común y natural de cambiar de plano en directo, ya que imita
el montaje cinematográfico y resulta imperceptible cuando el timing es correcto.

\subsection{Trans (transición suave)}

El botón \textbf{Trans} realiza una transición animada entre el estado actual y el
nuevo estado de composite o fuente, interpolando los parámetros geométricos y de
transparencia durante un tiempo configurable (por defecto 750~ms, definido en la
sección \texttt{[transitions]} del fichero INI). La transición se implementa
modificando progresivamente los valores de los pads del compositor de GStreamer.

Voctomix 2.0 define múltiples trayectorias de transición según el estado de origen y
destino (por ejemplo, de \texttt{fs} a \texttt{pip} la transición pasa por el composite
intermedio \texttt{fs-pip} antes de llegar a \texttt{pip}), garantizando que la
animación sea visualmente coherente independientemente del cambio solicitado.

\subsection{Señal de programa final}

La señal de programa final es la mezcla procesada por voctocore una vez aplicados el
composite activo y los overlays. Esta señal se emite simultáneamente en dos puertos:

\begin{itemize}
  \item \textbf{Puerto 11000}: salida de la mezcla en bruto (vídeo sin comprimir en
  I420), destinada a herramientas de grabación o encoders externos como FFmpeg.
  \item \textbf{Puerto 15000}: salida de streaming, que pasa previamente por el
  stream blanker. Cuando el estado es LIVE, esta salida es idéntica al puerto 11000;
  cuando el estado es PAUSE o NOSTREAM, emite el vídeo alternativo correspondiente.
\end{itemize}

El encoder de salida (tipicamente FFmpeg) se conecta al puerto 15000 y produce el
flujo H.264/AAC que se ingesta en la plataforma de streaming (RTMP/SRT hacia YouTube
Live, Twitch u otras plataformas).

% --- Figura sugerida ---
% \begin{figure}[H]
%   \centering
%   \includegraphics[width=0.7\textwidth]{figuras/botones_transicion.png}
%   \caption{Botones Cut y Trans en la barra de herramientas de voctogui.}
%   \label{fig:botones_transicion}
% \end{figure}
